                        %%%%%   PREAMBLE  %%%%%
%%%%%%%%%%%%%%%%%%%%%%%%%%%%%%%%%%%%%%%%%%%%%%%%%%%%%%%%%%%%%
%% HEADER
%%%%%%%%%%%%%%%%%%%%%%%%%%%%%%%%%%%%%%%%%%%%%%%%%%%%%%%%%%%%%

% \documentclass[12p,spanish,a4paper,oneside]{article}%Paquete de idioma español, problema con la separación de palabras
% Alternative Options:
%	Paper Size: a4paper / a5paper / b5paper / letterpaper / legalpaper / executivepaper
% Duplex: oneside / twoside
% Base Font Size: 10pt / 11pt / 12pt
\documentclass{beamer}
\usetheme{Darmstadt}
\usecolortheme{sidebartab}
\usepackage[sort]{natbib}

\usepackage{authblk}

%% Language %%%%%%%%%%%%%%%%%%%%%%%%%%%%%%%%%%%%%%%%%%%%%%%%%
\usepackage[spanish,activeacute]{babel} % english francais, polish, spanish, ...
\usepackage[T1]{fontenc}
%\usepackage[utf8]{inputenc} % Para uso en ingles
\usepackage[utf8]{inputenc}	%Para poder ingresar acentos y demás directamente
\usepackage{verbatim}%Paquete para poder comentar oraciones completas
% Por problemacon el idioma en las referencias ??
%\renewcommand\refname{{\large Referencias}} 
%\newcommand\VRule[1][\arrayrulewidth]{\vrule width #1}

\usepackage{lmodern} %Type1-font for non-english texts and characters

%% Packages for Graphics & Figures %%%%%%%%%%%%%%%%%%%%%%%%%%
\usepackage{graphicx} %%For loading graphic files
\usepackage{wrapfig} %% For text next to graphics
\usepackage{subfigure} %% Subfiguras
\usepackage[small,bf]{caption} %% Cambio de formato de los captions
%\usepackage{subfig} %%Subfigures inside a figure
%\usepackage{tikz} %%Generate vector graphics from within LaTeX
\usepackage{epsfig}

%% Please note:
%% Images can be included using \includegraphics{filename}
%% resp. using the dialog in the Insert menu.
%% 
%% The mode "LaTeX => PDF" allows the following formats:
%%   .jpg  .png  .pdf  .mps
%% 
%% The modes "LaTeX => DVI", "LaTeX => PS" und "LaTeX => PS => PDF"
%% allow the following formats:
%%   .eps  .ps  .bmp  .pict  .pntg


%% Line Spacing %%%%%%%%%%%%%%%%%%%%%%%%%%%%%%%%%%%%%%%%%%%%%
%\usepackage{setspace}
%\singlespacing        %% 1-spacing (default)
%\onehalfspacing       %% 1,5-spacing
%\doublespacing        %% 2-spacing


%% Other Packages %%%%%%%%%%%%%%%%%%%%%%%%%%%%%%%%%%%%%%%%%%%
%\usepackage{a4wide} %%Smaller margins = more text per page.
%\usepackage{fancyhdr} %%Fancy headings
%\usepackage{longtable} %%For tables, that exceed one page
%\usepackage{authblk} % For multi author affiliation support.
\newcommand\Mark[1]{\textsuperscript#1}

%% Math Packages %%%%%%%%%%%%%%%%%%%%%%%%%%%%%%%%%%%%%%%%%%%%
\usepackage{amsfonts}
\usepackage{amssymb,amsmath}
\usepackage{wasysym}
\usepackage{arcs} % Para poder escribir arcos... increible.
%\usepackage[amsthm,hyperref,thref]{ntheorem}


\usepackage{fancyhdr}
\renewcommand{\headrulewidth}{1pt}  %  línea horiz. bajo el encabezado
\renewcommand{\footrulewidth}{1pt}  %  línea horiz. sobre el pie



% \usepackage{makeidx}

%\usepackage{times}
%\usepackage[OT1,T1]{fontenc}

\usepackage{palatino}
% palatino: una fuente excelente, pero si sale
% mal en el .pdf descomentar las dos líneas anterires


%% Reseteo de Margenes
%\usepackage{vmargin}
%\setpapersize{A4}
%\setmarginsrb{2.5cm}{2cm}{1.5cm}{2cm}%
%{1cm}{1cm}%
%{1cm}{1cm}

%% Mejora de visualización
\usepackage{color}       % Para el color de las fuentes
\usepackage{lettrine}    % Letra capital
\usepackage{enumerate}   % Entorno enumerate mejorado
\usepackage{mdwlist}     % Entorno descripcion mejorado
\usepackage{hyperref}
\usepackage{datetime}	
\usepackage{listings}	 % Mejora de entorno para mostrar códigos
%\usepackage{paralist}	  
%\usepackage[framed,numbered,autolinebreaks,useliterate]{mcode}
%\usepackage[autolinebreaks,useliterate]{mcode}% Paquete de matlab para poder agregar el codigo
  

%\parskip 0.15cm          %  Queda mejor pero mas paginas


%\usepackage{appendix}	% Usando paquete de apendices

\usepackage{array,booktabs}
\usepackage{multirow}

% Para permitir un nuevo nivel de sectionado
\setcounter{secnumdepth}{5}

\usepackage{float}
\usepackage[toc]{appendix}


% CUSTOM COMMANDS 
% TODO NOTES
\usepackage{todonotes} % PARA AGREGAR NOTAS 'TODO' DENTRO DEL PDF
\newcommand{\todosection}[0]{\todo[inline]{Desarrollar Sección \thesection}}
\newcommand{\todosubsection}[0]{\todo[inline]{Desarrollar Subsección \thesubsection}}
\newcommand{\pendiente}[1]{\todo[inline]{\textbf{Pendiente:} #1}}
\newcommand{\pendpres}[1]{\todo[inline, size=\tiny]{{\textbf{Pendiente:} #1}}}

% CUSTOM COMMANDS
\newcounter{designlayer}
\newcommand{\designlayer}[2]{
	\refstepcounter{designlayer}\label{#2}
	\subsection{Etapa \thedesignlayer: #1}
}

\newcounter{requirement}
\newcommand{\requirement}[1]{
	\refstepcounter{requirement}\smallskip
	\textbf{Requerimiento \therequirement: }{{#1}.\smallskip}
}

\newcounter{activity}
\newcommand{\activity}[2]{
	\refstepcounter{activity}\smallskip \label{ac:#1}
	\textbf{Actividad \theactivity~(#1): }{#2.\smallskip}
}
\newcommand{\acref}[1]{
	{\textbf{A\ref{ac:#1} (#1)}}
}

\newcounter{component}
\newcommand{\component}[2]{
	\refstepcounter{component}\smallskip \label{cp:#1}
	\textbf{Componente \thecomponent~(#1): }{{#2}.\smallskip}
}
\newcommand{\cpref}[1]{
	{\textbf{C\ref{cp:#1} (#1)}}
}